\chapter{Massive Populations and Communication Degradation}
\label{sec:population}
% ===================================================================

\section{The Initial Condition}

We consider an \emph{initial condition} consisting of a population
$\pop_0$ of agents in an interactive GSLT $\GSLT$, with
$|\pop_0| \approx 10^{80}$ --- a number comparable to the number of
quarks in the observable universe.

\begin{remark}
The choice of scale is not arbitrary.
We want a population large enough that every cluster of agents is
itself a large enough system to implement a complex agent.
The quark-scale population ensures this at every level of the
hierarchy we are about to construct.
\end{remark}

\section{Communication Degradation}

\begin{definition}[Message Fidelity]
Let $A, B \in \pop$ be two agents and let $\gamma$ be a rewrite
path representing the transmission of a message from $A$ to $B$.
The \emph{fidelity} of the transmission is $|W(\gamma)|^2$, the
squared modulus of the path amplitude defined in Part~I.
A transmission is \emph{coherent} if its fidelity exceeds a threshold
$\epsilon > 0$.
\end{definition}

\begin{definition}[Communication Radius]
The \emph{communication radius} of an agent $A \in \pop$ is the set
\[
  B_\epsilon(A) = \{ B \in \pop \mid
  \exists \text{ coherent path } \gamma : A \to B \}
\]
of agents reachable from $A$ with fidelity above $\epsilon$.
\end{definition}

\begin{conjecture}[Degradation at Scale]
\label{conj:degradation}
For any interactive GSLT $\GSLT$ with bounded step amplitudes,
and for any $\epsilon > 0$, there exists a scale $N_\epsilon$ such
that for any population $\pop$ with $|\pop| > N_\epsilon$:
\[
  \exists A \in \pop \text{ such that }
  B_\epsilon(A) \subsetneq \pop.
\]
That is, no single agent can communicate coherently with the entire
population.
\end{conjecture}

\begin{remark}
The intuition is that message fidelity degrades along each
transmission hop by a factor bounded away from $1$.
Over a path of length $n$ the fidelity is at most $(1-\delta)^n$
for some $\delta > 0$.
For a population of diameter $d$ (minimum path length between
the most distant agents), fidelity falls below $\epsilon$ when
$d > \log(1/\epsilon) / \log(1/(1-\delta))$.
In a population of size $10^{80}$ the diameter is enormous,
and no $\epsilon$-coherent path can span it.
Making this precise requires specifying the communication geometry
of the GSLT --- a gap noted in Section~\ref{sec:gaps}.
\end{remark}

\begin{remark}[No Error-Correction Escape]
One might hope that error-correction protocols could restore
global coherence.
We argue this hope fails at sufficient scale.
Any error-correction protocol is itself a computation in $\GSLT$,
consuming account resources and communicating over the same
degraded channels.
At population scales beyond $N_\epsilon$, the overhead of
error-correction exceeds the bandwidth of the channels it is trying
to protect.
The argument is self-referential: error-correction cannot bootstrap
coherence it does not already have.
This argument is currently informal; making it precise is a
priority for future work.
\end{remark}

% ===================================================================
